\documentclass[aps,prd,reprint,onecolumn,superscriptaddress,nofootinbib]{revtex4-2}
\usepackage{amsmath,amssymb,bm,graphicx,booktabs,array,hyperref}
\hypersetup{colorlinks=true,linkcolor=blue,citecolor=blue,urlcolor=blue}
\newcommand{\dd}{\mathrm d}
\newcommand{\PiH}{\mathsf P}
\begin{document}
\title{Cross-Scale Geometric Representation of Classical Inspirals:\\Four-Dimensional Helical Embeddings and Mass-Substitution Fingerprints\\{\small SANDBOXED RESEARCH MANUSCRIPT}}
\author{Nathan McKnight}
\affiliation{Independent Researcher, Columbus, Ohio, USA}
\date{October 11, 2026}
\begin{abstract}
We investigate a representation-level comparison of established inspiral equations without altering their interactions or adding dynamical terms. The leading-order gravitational-wave-driven circular-binary inspiral and the adiabatic, nonrelativistic Coulomb--Larmor inspiral are each evaluated under two orbiting-mass substitutions, using an electron mass and an Earth mass while holding the remaining parameters fixed within each equation family. The resulting planar trajectories are embedded as curves in four-dimensional Euclidean coordinate space, $X(t)=(ct,r\cos\phi,r\sin\phi,0)$. Their common kinematic parameterization is $\dot r=-ar^{-n}$ and $\dot\phi=br^{-3/2}$, with $n=3$ and $n=2$, respectively. An interaction-neutral pitch observable, $\PiH=|\dot r|/(r\dot\phi)$, satisfies $\PiH_G=(64/5)\nu\beta^5$ for the gravitational model and $\PiH_E=(4/3)|q/Q|\beta^3$ for the electromagnetic model, where $\beta=r\dot\phi/c$ and $\nu$ is the symmetric mass ratio. We establish exact mappings between normalized winding profiles, distinguish them from four-dimensional similarities, and derive logarithmic scaling fingerprints. A damped planar oscillator provides an independent comparator. The results define a reproducible protocol for cross-scale geometrical inspection of known solutions; they do not identify the physical interactions or assert a new dynamical unification.
\end{abstract}
\maketitle

\section{Introduction}
A comparison between physical theories need not begin by equating their interactions. One can instead compute the behavior prescribed by each theory, express the behavior in a common geometrical language, suppress the labels identifying its origin, and compare minimal pairs of geometrical objects. This procedure is motivated by the mapping-first strategy developed in the Scalar--Angular--Torsion (SAT) research program. Here we isolate a narrow and independently reproducible mathematical test: the comparison of inward-spiraling trajectories from already established classical equations.

Our additional operation is deliberately unsophisticated: replace the orbiting mass in each canonical equation by each of two values differing by approximately 55 orders of magnitude, leaving the equation and its other specified parameters unchanged. The purpose is to expose the response already encoded in the equations rather than to insert an interpolating or compensating scaling law. The cross-scale cases are formal mathematical evaluations: in particular, the classical Coulomb model is not a model of a stationary quantum atomic electron, and an Earth-mass object assigned one elementary charge and an atomic orbital radius is not a realizable macroscopic charged body.

The contribution is methodological and analytic: (i) a common four-dimensional helical embedding, (ii) an exact algebraic comparison of pitch--winding laws, (iii) transformations of normalized angular and temporal profiles, (iv) a distinction between geometric similarities and mere curve correspondences, and (v) a general logarithmic response fingerprint checked against a third classical spiral model. The paper intentionally excludes hypotheses about an underlying common force or worldtube microstructure.

\section{Canonical equations and controlled substitutions}
\subsection{Gravitational radiation reaction}
For a quasi-circular binary with component masses $M$ and $m$, separation $r$, and leading quadrupole gravitational-wave dissipation, the standard adiabatic evolution is \cite{peters1964,maggiore2008}
\begin{align}
 \dot r_G&=-\frac{64G^3}{5c^5}\frac{Mm(M+m)}{r^3}\equiv-\frac{a_G}{r^3},\\
 \dot\phi_G&=\sqrt{\frac{G(M+m)}{r^3}}\equiv\frac{b_G}{r^{3/2}},\label{eq:G}
\end{align}
with $a_G=64G^3Mm(M+m)/(5c^5)$ and $b_G=[G(M+m)]^{1/2}$. We hold $M=1.98847\times10^{30}\,\mathrm{kg}$ and $r_{0G}=1\,\mathrm{AU}$ fixed and substitute two different orbiting masses $m$. The radial equation is a leading-order, adiabatic approximation, not a full numerical-relativity trajectory.

\subsection{Classical electromagnetic radiation reaction}
For an orbiting charge $q$ of inertial mass $m$ attracted to a fixed central charge $Q$ of opposite sign, let $K=|qQ|/(4\pi\epsilon_0)$. The circular Coulomb balance and nonrelativistic Larmor power are \cite{jackson1999}
\begin{align}
 \dot\phi_E&=\sqrt{\frac{K}{mr^3}}, &
 P_L&=\frac{q^2 a_{\mathrm{rad}}^2}{6\pi\epsilon_0c^3}, &
 a_{\mathrm{rad}}&=\frac{K}{mr^2}.
\end{align}
Using the circular orbital energy $E=-K/(2r)$ and $\dot E=-P_L$ yields
\begin{equation}
 \dot r_E=-\frac{q^2K}{3\pi\epsilon_0 c^3m^2r^2}\equiv-\frac{a_E}{r^2},\qquad
 \dot\phi_E=\frac{b_E}{r^{3/2}},\label{eq:E}
\end{equation}
where $a_E=q^2K/(3\pi\epsilon_0 c^3m^2)$ and $b_E=(K/m)^{1/2}$. We hold $|q|=|Q|=e$ and $r_{0E}=5.29177210903\times10^{-11}\,\mathrm m$ fixed. This is the formal adiabatic circular-orbit model, neglecting retardation beyond the adopted radiation-loss law, relativistic corrections, and quantum mechanics. Its applicability is conditional on slowly changing bound motion.

\subsection{Four-case design}
The only varied physical input \emph{within each equation family} is $m$: the electron mass $m_e=9.1093837139\times10^{-31}\,\mathrm{kg}$ or the Earth mass $m_\oplus=5.9722\times10^{24}\,\mathrm{kg}$. The ratio is $m_\oplus/m_e\simeq6.56\times10^{54}$. The gravitational and electromagnetic setups have different central sources and starting radii, inherited from their separate reference configurations; comparisons across the two families are therefore not single-variable controlled physical experiments. No values are fitted to optimize their agreement.

\begin{table}[h]
\caption{Original four-map numerical inputs and outputs. $\beta_0=r_0\dot\phi_0/c$ and $\PiH_0=|\dot r_0|/(r_0\dot\phi_0)$. $T_0$ and $N_0$ denote formal, extrapolated limits to $r=0$ and should not be interpreted as physical endpoint predictions.}
\begin{ruledtabular}
\begin{tabular}{lcccc}
Map and description & $\beta_0$ & $\PiH_0$ & $T_0$ (s) & $N_0$ (revolutions)\\
A: G, $m_e$ & $9.9352540\times10^{-5}$ & $5.676424\times10^{-80}$ & $2.21203\times10^{85}$ & $1.12152\times10^{78}$\\
B: G, $m_\oplus$ & $9.9352689\times10^{-5}$ & $3.721524\times10^{-25}$ & $3.37399\times10^{30}$ & $1.71064\times10^{23}$\\
C: E, $m_e$ & $7.2973526\times10^{-3}$ & $5.181252\times10^{-7}$ & $1.55618\times10^{-11}$ & $2.04783\times10^{5}$\\
D: E, $m_\oplus$ & $2.8499855\times10^{-30}$ & $3.086503\times10^{-89}$ & $6.68882\times10^{98}$ & $3.43765\times10^{87}$
\end{tabular}
\end{ruledtabular}
\end{table}

\section{Common four-dimensional geometric representation}
Each planar solution is assigned the Euclidean-coordinate embedding
\begin{equation}
 X(t)=\big(ct,\,r(t)\cos\phi(t),\,r(t)\sin\phi(t),\,0\big).\label{eq:X}
\end{equation}
This is a curve in $\mathbb R^4$ equipped, for the present geometrical comparison, with the Euclidean metric $\delta_{AB}$. The coordinate $x^0=ct$ is measured in meters. The embedding does \emph{not} replace the Lorentzian physical spacetime metric by a proposed dynamical Euclidean metric, and is not itself evidence of four-dimensional SAT worldtube dynamics. Only the centerlines are defined; no finite-radius core or cross-section is supplied.

Both examples have the shared evolution form
\begin{equation}
 \dot r=-ar^{-n},\qquad \dot\phi=br^{-3/2},\label{eq:common}
\end{equation}
with $n_G=3$ and $n_E=2$ and inherited $a,b>0$. Direct integration from $(r_0,\phi_0)$ gives
\begin{align}
 r(t)&=\left[r_0^{n+1}-(n+1)at\right]^{1/(n+1)},\label{eq:rsol}\\
 \phi(r)-\phi_0&=\frac{b}{a(n-\tfrac12)}\left[r_0^{n-1/2}-r^{n-1/2}\right].\label{eq:phisol}
\end{align}
The formal extrapolated timescale and winding are
\begin{equation}
 T_0=\frac{r_0^{n+1}}{(n+1)a},\qquad N_0=\frac{b\,r_0^{n-1/2}}{2\pi a(n-1/2)}.
\end{equation}
Equations (\ref{eq:rsol})--(\ref{eq:phisol}) are valid mathematical continuations only as far as their original approximations permit.

Define interaction-neutral local descriptors
\begin{equation}
 \beta\equiv\frac{r\dot\phi}{c},\qquad
 \PiH\equiv\left|\frac{\dot r}{r\dot\phi}\right|=\frac{a}{b}r^{1/2-n}.
 \label{eq:descriptors}
\end{equation}
The full Euclidean tangent angle relative to the $ct$ axis is $\tan\theta=\sqrt{\dot r^2+(r\dot\phi)^2}/c=\beta\sqrt{1+\PiH^2}$; $\beta$ alone records its azimuthal component. Coordinates and their dimensional conventions are held fixed for all four embeddings.

\section{Geometric response laws}
For the gravitational system, let $\nu=Mm/(M+m)^2$. Since $\beta_G=\sqrt{G(M+m)/(rc^2)}$, Eqs. (\ref{eq:G}) and (\ref{eq:descriptors}) imply the exact identity \emph{within the leading-order model}
\begin{equation}
 \boxed{\PiH_G=\frac{64}{5}\nu\beta_G^5.}\label{eq:pG}
\end{equation}
For the classical electromagnetic model, $\beta_E=\sqrt{K/(mrc^2)}$ gives
\begin{equation}
 \boxed{\PiH_E=\frac43\left|\frac qQ\right|\beta_E^3.}\label{eq:pE}
\end{equation}
Thus $\dd\ln\PiH/\dd\ln\beta=5$ and $3$ along each inspiral with its respective other parameters fixed. At an artificially matched value of $\beta$, the \emph{functions} obey
\begin{equation}
 \frac{\PiH_G}{\PiH_E}=\frac{48}{5}\frac{\nu}{|q/Q|}\beta^2.\label{eq:ratio}
\end{equation}
Equation (\ref{eq:ratio}) is a comparison of two model functions at a shared argument, not the measured ratio of any particular cross-scale pair. No universal scaling exponent follows from this identity alone.

\begin{figure}[t]
\includegraphics[width=.79\linewidth]{fig_pitch.pdf}
\caption{Analytic contraction-pitch profiles of the two orbital models and an independent damped-oscillator comparator. Each pitch is divided by its own initial value \emph{after} solving the equations. The slopes $-5/2$, $-3/2$, and $0$ are equation-derived, not fitted.}\label{fig:pitch}
\end{figure}

\section{Normalized correspondences and their limitations}
For $x=r/r_0$, the fraction of formal accumulated angular winding is
\begin{equation}
 f_{\phi}(x)=\frac{\phi(r)-\phi_0}{\phi(0)-\phi_0}=1-x^{n-1/2}.
\end{equation}
The gravitational and electromagnetic profiles are $1-x_G^{5/2}$ and $1-x_E^{3/2}$. Requiring equal \emph{fractions} of accumulated winding gives
\begin{equation}
 x_E=x_G^{5/3}.\label{eq:phaseMap}
\end{equation}
In contrast, equal fractions of formal elapsed time $u=t/T_0=1-x^{n+1}$ require
\begin{equation}
 x_E=x_G^{4/3}.\label{eq:timeMap}
\end{equation}
These are distinct transformations. Under the winding correspondence (\ref{eq:phaseMap}), the residual in normalized elapsed time is
\begin{equation}
 \Delta u=\big[1-x_G^5\big]-\big[1-x_G^4\big]=x_G^4(1-x_G),\qquad
 \max\Delta u=\frac{256}{3125}=0.08192
\end{equation}
at $x_G=4/5$. There is no single \emph{radial power substitution} aligning both independently normalized profiles. More general nonlinear mappings can map parameterized curves into one another, so this observation must not be described as a proof that no four-dimensional diffeomorphism exists.

\begin{figure}[t]
\includegraphics[width=.72\linewidth]{fig_phase.pdf}
\caption{Angular winding fractions. The post-solution nonlinear substitution $x_E=x_G^{5/3}$ overlays the normalized winding profiles exactly, but does not align their physical radial scales or actual total winding angles.}\label{fig:phase}
\end{figure}
\begin{figure}[t]
\includegraphics[width=.72\linewidth]{fig_time.pdf}
\caption{Difference in normalized elapsed time after matching the normalized angular winding fractions. The maximum residual is $0.08192$.}\label{fig:time}
\end{figure}

More explicitly, the phase-aligned correspondence has local factors
\begin{align}
 \frac{\dd r_E}{\dd r_G}&=\frac53\frac{r_{0E}}{r_{0G}}x_G^{2/3},\\
 \frac{\dd t_E}{\dd t_G}&=\frac54\frac{T_E}{T_G}x_G,\\
 \frac{\dd\phi_E}{\dd\phi_G}&=\frac{\Phi_E}{\Phi_G},
\end{align}
where $\Phi$ denotes total formal angular winding. These are coordinate- and normalization-dependent stretching factors, not a Euclidean similarity transformation.

\section{Intrinsic Euclidean differential geometry}
The embedded curve has Euclidean arclength
\begin{equation}
 \dd\sigma=\sqrt{c^2+\dot r^2+r^2\dot\phi^2}\,\dd t,\qquad
 \left|\frac{\dd\sigma}{\dd r}\right|=\sqrt{1+\frac{c^2r^{2n}}{a^2}+\frac{b^2r^{2n-1}}{a^2}}.\label{eq:ds}
\end{equation}
Because the trajectory lies in the three-dimensional subspace $(ct,x,y)$, its usual Frenet curvature and signed torsion can be calculated with $V=\dd X/\dd t$, $A=\dd^2X/\dd t^2$, and $J=\dd^3X/\dd t^3$:
\begin{equation}
 \kappa=\frac{\|V\times A\|}{\|V\|^3},\qquad
 \tau=\frac{\det(V,A,J)}{\|V\times A\|^2}.\label{eq:frenet}
\end{equation}
The cross product here acts only in that explicitly specified three-dimensional subspace. The third generalized Frenet curvature vanishes identically for these planar-spatial embeddings.

Under a positive uniform Euclidean dilation $X\mapsto\lambda X$, $\kappa\mapsto\kappa/\lambda$, $\tau\mapsto\tau/\lambda$, and $\sigma\mapsto\lambda\sigma$. Two dimensionless similarity invariants, where $\tau\ne0$, are
\begin{equation}
 h=\frac{\kappa}{|\tau|},\qquad
 \mathcal S=\frac{1}{\kappa}\frac{\dd h}{\dd\sigma}.\label{eq:intrinsic}
\end{equation}
For a slowly degrading, nonrelativistic helix $\PiH\ll1$ and $\beta\ll1$, the circular-helix expressions yield $h\simeq\beta$ and $\kappa\simeq\beta^2/r$. Since $\beta\propto r^{-1/2}$ and $\dd r/\dd\sigma\simeq-\beta\PiH$, differentiation gives
\begin{equation}
 \mathcal S\simeq\frac{\PiH}{2},\qquad
 \mathcal S_G\simeq\frac{32}{5}\nu h^5,\qquad
 \mathcal S_E\simeq\frac23\left|\frac qQ\right|h^3.\label{eq:S}
\end{equation}
These last formulas are \emph{asymptotic}, not exact identities of degrading helices. Their agreement with the simpler pitch laws at small degradation does not demonstrate a global isometry. The coefficient-dependent geometric relationships are distinguished from dimensionless invariants themselves.

\section{Logarithmic scaling fingerprints}
For any contracting planar curve represented by
\begin{equation}
 \dot r=-F(r,m),\qquad \dot\phi=\Omega(r,m),\qquad F,\Omega>0,
\end{equation}
define the coordinate-specified response fingerprint
\begin{equation}
 \bm{\mathfrak f}=\left(\frac{\partial\ln F}{\partial\ln r},\frac{\partial\ln F}{\partial\ln m},\frac{\partial\ln\Omega}{\partial\ln r},\frac{\partial\ln\Omega}{\partial\ln m}\right).\label{eq:fp}
\end{equation}
Partial derivatives hold the independent source parameters fixed. For our established equations,
\begin{align}
 \bm{\mathfrak f}_G&=\left(-3,1+\frac m{M+m},-\frac32,\frac{m}{2(M+m)}\right),\\
 \bm{\mathfrak f}_E&=\left(-2,-2,-\frac32,-\frac12\right).\label{eq:fps}
\end{align}
For any such representation $\PiH=F/(r\Omega)$, and therefore
\begin{equation}
 \frac{\partial\ln\PiH}{\partial\ln r}=\frac{\partial\ln F}{\partial\ln r}-1-\frac{\partial\ln\Omega}{\partial\ln r}.\label{eq:identity}
\end{equation}
Equation (\ref{eq:identity}) is a kinematic identity; the different numerical slopes of $-5/2$ and $-3/2$ are consequences of the input models. The fingerprint is not an intrinsic geometric invariant: it specifies responses with respect to named physical parameters and coordinates. Recovering a mass-response derivative requires \emph{multiple} mass-substituted solutions, not one anonymous curve.

\begin{figure}[t]
\includegraphics[width=.75\linewidth]{fig_fingerprint.pdf}
\caption{Logarithmic fingerprints of the two canonical orbit models evaluated at electron mass. The four entries are coordinate- and parameter-specified partial sensitivities, not Euclidean invariants.}\label{fig:fp}
\end{figure}

\section{Independent comparator: a damped circular oscillator}
As a check that the representation is not specific to inverse-square orbital attraction, consider the planar linear oscillator
\begin{equation}
 m\ddot{\bm x}+\gamma\dot{\bm x}+k\bm x=0,
\end{equation}
with fixed positive $k,\gamma$ and the underdamped condition $\gamma^2<4km$. A circularly polarized exact solution is
\begin{equation}
 r(t)=r_0e^{-\gamma t/(2m)},\qquad
 \phi(t)=\omega_dt,\qquad
 \omega_d=\sqrt{\frac{k}{m}-\frac{\gamma^2}{4m^2}}.
\end{equation}
Hence $F_D=\gamma r/(2m)$, $\Omega_D=\omega_d$, and
\begin{equation}
 \bm{\mathfrak f}_D=\left(1,-1,0,-\frac12+\frac{\gamma^2}{8m^2\omega_d^2}\right),\qquad
 \frac{\partial\ln\PiH_D}{\partial\ln r}=0.
\end{equation}
The oscillator generates a contracting helical embedding, but its pitch remains constant as a function of radius. This counterexample shows that common helical \emph{representation} does not force common scaling exponents. Any two-mass oscillator demonstration must also check that the system remains underdamped at both substituted masses.

\section{Discussion and scope}
The analysis identifies two levels of commonality. First, a broad class of planar degrading motions admits the kinematic embedding (\ref{eq:X}), with measurable winding, radial contraction, and Euclidean Frenet properties. Second, equation-specific parameter dependencies persist within that representation. The gravitational and electromagnetic examples share $\dot\phi\propto r^{-3/2}$ but differ in radial contraction, giving the fifth- and third-power pitch relationships. A third spiral violates even the shared angular-radius exponent.

The nonlinear substitution (\ref{eq:phaseMap}) is a precise correspondence between \emph{normalized functions}; it is not a physical conversion between interactions, not a coordinate-independent characterization of scale, and not a derivation of an interaction-unifying law. The use of Euclidean $\delta_{AB}$ here is a choice of geometric description for a curve. Testing a physical SAT/H(s)H worldtube hypothesis would require additional structure and new empirical comparisons; neither is part of this paper.

A formal mass substitution can also leave the domain of validity of the underlying approximation. Our very small and very large formal inspiral times and total winding counts, especially in the electromagnetic cross-scale case, are therefore bookkeeping properties of extended equations rather than experimental forecasts. The paper does not claim a quantum electron executes a classical radiatively collapsing orbit, or that the quadrupole or Larmor approximations survive all the way to $r=0$.

The operational advantage of the present construction is reproducibility. Within an equation family, the substitutions are fixed in advance and the geometric mapping contains no tunable coefficients. The normalized response laws and their failures of simple similarity can be audited independently, and the same procedure can be tested against additional canonical systems with known trajectories.

\section{Conclusion}
Na\"ive orbiting-mass substitution in two established inspiral approximations yields four formally computable degrading trajectories. Their four-dimensional Euclidean helical embeddings share a common kinematic form while retaining distinct geometric scaling behavior. Specifically, gravitational and electromagnetic orbital-decay models give $\PiH=(64/5)\nu\beta^5$ and $\PiH=(4/3)|q/Q|\beta^3$; their normalized phase histories correspond under $x_E=x_G^{5/3}$, whereas equal normalized elapsed time requires $x_E=x_G^{4/3}$. The two mappings are not ordinary four-dimensional similarities. A damped oscillator confirms that helical geometry alone does not dictate the scaling fingerprint. The generalizable contribution is a parameter-explicit procedure for comparing known dynamics geometrically before inferring any scaling relationship. No new physical interactions or dynamical assumptions are required.

\begin{acknowledgments}
The author acknowledges the extended SAT/H(s)H research archive and AI-assisted symbolic and numerical drafting. Derivations, conventions, and claims in this manuscript are intended to be independently auditable using the accompanying code. This document is a sandboxed research manuscript, not a claim of peer review.
\end{acknowledgments}
\appendix
\section{Reproducibility and dimensional conventions}
The Python companion \texttt{reproduce.py} uses SciPy constants for $G,c,\epsilon_0,e,m_e$ and the fixed values of $M,m_\oplus,r_{0G},r_{0E}$ specified above. It checks Eqs. (\ref{eq:pG}), (\ref{eq:pE}), the normalized winding correspondence, and the oscillator differential equation symbolically with SymPy. It exports \texttt{four\_maps.csv} and the plotted figures. The dimension of $a$ is $\mathrm m^{n+1}\mathrm s^{-1}$; that of $b$ is $\mathrm m^{3/2}\mathrm s^{-1}$. All displayed exponents apply to these exact input equations and stated variable-holding conventions.

The formal endpoint $r=0$ is used only to normalize analytic expressions; it does not indicate the physical validity of the source equations there. Numerical floating-point samples of very large total winding counts cannot be reconstructed reliably by reducing enormous phases modulo $2\pi$; the analyses here use analytic winding differences and local derivatives instead.

\begin{thebibliography}{9}
\bibitem{peters1964} P. C. Peters, ``Gravitational Radiation and the Motion of Two Point Masses,'' \textit{Physical Review} \textbf{136}, B1224--B1232 (1964), doi:10.1103/PhysRev.136.B1224.
\bibitem{jackson1999} J. D. Jackson, \textit{Classical Electrodynamics}, 3rd ed. (Wiley, New York, 1999), chapters on accelerated charges and radiation.
\bibitem{maggiore2008} M. Maggiore, \textit{Gravitational Waves, Volume 1: Theory and Experiments} (Oxford University Press, Oxford, 2008), sections on quadrupole emission and compact binaries.
\end{thebibliography}
\end{document}